\section{Flags on a point blowup}
\label{sec:auxiliary-blowup-flag}

We now remove the assumption that $X$ itself carries a complete compact flag. In dimension at least two, the exceptional divisor of a point blowup carries a complete projective flag. Hartogs extension preserves holomorphic sections, so the growth formula on the modification counts the original section spaces. On a curve, a point already completes the flag.

\begin{lemma}\label{lem:point-blowup-section-descent}
Every connected compact complex manifold $X$ of dimension $n\geq1$ admits a smooth compact modification $\pi:\widetilde X\to X$ with a complete smooth compact flag. It may be chosen so that, for every holomorphic line bundle $L$ and every $m\geq0$, pullback gives
\[
 H^0(X,L^m)\simeq H^0(\widetilde X,\pi^*L^m).
\]
These identifications respect products.
\end{lemma}
\begin{proof}
If $n=1$, take $\pi$ to be the identity and choose any point as the last member of the flag. If $n\geq2$, blow up a point $p\in X$. The usual local construction of the blowup of $\mathbb C^n$ at the origin shows that $\widetilde X$ is smooth and its exceptional divisor is $E\simeq\mathbb P^{n-1}$. Choose a nested linear flag in $E$. Together with $\widetilde X$ it gives
\[
 \widetilde X\supset E\supset\mathbb P^{n-2}
       \supset\cdots\supset\mathbb P^0.
\]
Every member is smooth and compact, and is a prime divisor in the preceding member. The modification is an isomorphism off $E$.

Let $s$ be a section of $\pi^*L^m$. Off $E$, it gives a section of $L^m$ on $X\setminus\{p\}$. In a trivialization of $L^m$ on a coordinate ball at $p$, its coefficient is holomorphic on the punctured ball. The classical Hartogs extension theorem in dimension $n\geq2$ extends that coefficient holomorphically across $p$. It glues with the section away from $p$. Its pullback agrees with $s$ off $E$, and hence everywhere by the identity theorem. This proves surjectivity; injectivity and product compatibility follow from equality on the same dense open set. The identity-map case is immediate.
\end{proof}

On the modification in \cref{lem:point-blowup-section-descent}, apply the flag construction to $\widetilde W_m=\pi^*W_m\subseteq H^0(\widetilde X,\pi^*L^m)$. The resulting body's shape may depend on the point and flag; its dimension and normalized count are determined by the original section algebra.

\begin{theorem}\label{thm:all-compact-section-count}
Assume that some positive-degree piece of $W_\bullet$ is nonzero. There are an integer $0\leq d\leq n$, a degree period $e\geq1$, and a number $V(W_\bullet)>0$ such that
\[
 \lim_{q\to\infty}\frac{\dim W_{eq}}{(eq)^d}=V(W_\bullet),
 \qquad W_m=0\quad\hbox{if }e\nmid m.
\]
Here $d$ is the transcendence degree of the section-ratio field $K(W_\bullet)$, $e$ is the greatest common divisor of the nonzero degrees, and $V(W_\bullet)$ is the intrinsic lattice volume of every auxiliary body constructed above. In particular these three quantities are independent of the auxiliary modification and flag.
\end{theorem}
\begin{proof}
Pullback preserves the graded dimensions and the nonzero degrees. Apply \cref{thm:gauduchon-linear-flag-bound,thm:intrinsic-flag-count} on $\widetilde X$ to obtain the displayed limit and positive lattice volume. Pullback also identifies the section-ratio fields: their elements are the corresponding ratios, and injectivity follows on the dense isomorphism locus. Thus \cref{prop:flag-dimension-section-field} identifies $d$ with the original field's transcendence degree. The unchanged nonzero degrees determine $e$, and the limit of the unchanged dimensions determines $V(W_\bullet)$. All three are independent of the auxiliary choices.
\end{proof}

\Needspace{16\baselineskip}
\begin{corollary}\label{cor:non-moishezon-section-vanishing}
For every holomorphic line bundle $L$ on a connected compact complex manifold $X$ of dimension $n$,
\[
 h^0(X,L^m)=O(m^n).
\]
If $X$ is not Moishezon, then, for each $L$, either all its positive-degree section spaces vanish, or there is an integer $d<n$ for which
\[
 h^0(X,L^m)=O(m^d).
\]
In particular $m^{-n}h^0(X,L^m)\to0$. These conclusions apply when $X$ fails the bounded mass property.
\end{corollary}
\begin{proof}
Apply the theorem to the full section ring, taking the zero-ring-in-positive-degrees case separately. The periodic count gives $O(m^d)$ for all degrees, with $d\leq n$. If $d=n$, the section-ratio field supplies $n$ algebraically independent meromorphic functions on $X$. By Moishezon's theorem, $X$ admits a projective modification and is Moishezon. Thus a non-Moishezon manifold has $d<n$. Finally, a Moishezon manifold has the bounded mass property by \cref{lem:kahler-modification-bounded-mass}. Its contrapositive proves the last assertion.
\end{proof}

The same upper bounds apply to multiplier-ideal section spaces by their inclusion in $H^0(X,L^m)$. Identifying weighted spaces on a modification additionally requires the Jacobian correction; \cref{lem:adjoint-ideal-transform} gives that formula with the relative canonical bundle.
